\chapter{Conclusions}
The established radiation-monitor market remains concentrated in LEO, MEO and GEO, where operational requirements are comparatively mature. Silicon telescopes and multi-sensor suites remain the most defensible solution when species separation, wide energy coverage and calibration fidelity are required. Pixel detectors have expanded the low-SWaP design space, particularly for event imaging, distributed monitoring and small spacecraft, but they do not automatically replace thicker telescope or calorimeter architectures.

The emerging cislunar market changes the balance of requirements rather than the underlying detector physics. Instruments must operate outside continuous geomagnetic shielding, may need to distinguish incident from albedo populations, and increasingly combine charged-particle, neutron and gamma-ray measurements within restricted accommodation budgets. The strongest opportunity here being pursued is not a single universal sensor, but a modular, cross-calibrated family spanning operational dosimetry, compact spectrometry and higher-fidelity science measurements. However missions like ARGONAUT and NovaMoon and similar offer opportunities to be on service providing platforms for those exploring the Moon.

%The next stage should add verified 2023--2026 status updates for every heritage instrument family, numerical mass and power fields suitable for plotting, response-matrix and calibration references, and a formal technology-readiness assessment based on complete instrument configurations rather than detector components alone.
